% sections/03_taxonomy.tex
% §3: Inconsistency Checking Framework (IC Framework)
% Restructured to follow IC difference-matrix approach (Liao et al., ACSAC 2020)

\section{Inconsistency Checking Framework}
\label{sec:taxonomy}

Inspired by the inconsistency checking methodology of~\cite{liao2020measuring},
we formalize capability-description mismatches for LLM agent skills
via a \emph{difference matrix} that captures both directions of divergence
between what a skill \emph{claims} (user-visible) and what it \emph{does} (actual).
Figure~\ref{fig:methodology} illustrates the extraction and classification pipeline.

\begin{figure}[t]
\centering
\includegraphics[width=\columnwidth]{figures/fig_methodology_overview.pdf}
\caption{Methodology overview: each SKILL.md is decomposed into description (yielding $\mathcal{C}_{\text{vis}}$) and code blocks (yielding $\mathcal{C}_{\text{act}}$ via LSCPS). The difference matrix produces four quadrants (AD, UD, FC, TN), from which IC types and quantitative metrics are derived.}
\label{fig:methodology}
\end{figure}

% ── Notation Table ──────────────────────────────────────────────────────────
\begin{table}[h]
\centering
\caption{Notation. All metrics are formally defined in \S\ref{sec:taxonomy}.}
\label{tab:notation}
\small
\begin{tabular}{@{}lp{6.2cm}@{}}
\toprule
\textbf{Symbol} & \textbf{Definition} \\
\midrule
$\mathcal{C}_{\text{vis}}(s)$  & User-visible capability set: mentions in marketplace description \\
$\mathcal{C}_{\text{act}}(s)$  & Actual capability set: code-detected capabilities across all skill artifacts \\
$\mathcal{F}$                  & Universe of 14 capability families (Table~\ref{tab:lscps_kw}) \\
\midrule
$\text{AD}(s)$ & Accurate Disclosure: $|\mathcal{C}_{\text{act}} \cap \mathcal{C}_{\text{vis}}|$ \\
$\text{UD}(s)$ & Undisclosed: $|\mathcal{C}_{\text{act}} \setminus \mathcal{C}_{\text{vis}}|$ (hidden) \\
$\text{FC}(s)$ & False Claim: $|\mathcal{C}_{\text{vis}} \setminus \mathcal{C}_{\text{act}}|$ (over-claimed) \\
$\text{TN}(s)$ & True Negative: $|\mathcal{F} \setminus (\mathcal{C}_{\text{act}} \cup \mathcal{C}_{\text{vis}})|$ \\
\midrule
$\text{CR}(s)$  & Capability Recall (Eq.~\ref{eq:cr}) \\
$\text{CP}(s)$  & Description Precision (Eq.~\ref{eq:cp}) \\
$\text{CDI}(s)$ & Capability Discrepancy Index (Eq.~\ref{eq:cdi}) \\
$\text{RWM}(s)$ & Risk-Weighted Mismatch (Eq.~\ref{eq:rwm}) \\
$\text{TCE}$    & Temporal Capability Expansion (Eq.~\ref{eq:tce}) \\
$\text{SDS}$    & Semantic Disclosure Score (changelog coverage) \\
\bottomrule
\end{tabular}
\end{table}

% ── Information Sources ──────────────────────────────────────────────────────
\subsection{User-Visible vs.\ Actual Capability Sets}

A LLM agent skill exposes information through two distinct channels:

\paragraph{User-visible layer ($\mathcal{C}_{\text{vis}}$).}
The marketplace description is the \emph{only} information a user
reads before installation. From this natural-language text,
we extract the set of capability families mentioned,
using a keyword dictionary derived from the LSCPS pattern bank (\S\ref{sec:static}).
This mirrors how~\cite{liao2020measuring} extracted data practices
from voice-app descriptions.

\paragraph{Actual capability layer ($\mathcal{C}_{\text{act}}$).}
The SKILL.md file contains the skill's executable specification,
invisible to end users. In Phase~1 (this paper),
$\mathcal{C}_{\text{act}} = \mathcal{C}_{\text{code}}$,
derived via Static Capability Screening (SCS) over all code blocks.
Phase~2 (future work) extends this to include natural-language
instruction grants ($\mathcal{C}_{\text{instr}}$) and
prose capability descriptions ($\mathcal{C}_{\text{text}}$) within SKILL.md:
\[
  \mathcal{C}_{\text{act}} \supseteq \mathcal{C}_{\text{code}}
  \;\cup\; \mathcal{C}_{\text{instr}} \;\cup\; \mathcal{C}_{\text{text}}
\]
Restricting to $\mathcal{C}_{\text{code}}$ makes our measurements
\emph{conservative}: the true mismatch is at least as large as reported.

% ── Difference Matrix ────────────────────────────────────────────────────────
\subsection{The Inconsistency Difference Matrix}

\begin{definition}[Capability Difference Matrix]
\label{def:dm}
For skill $s$ with capability sets $\mathcal{C}_{\text{act}}(s)$ (actual) and
$\mathcal{C}_{\text{vis}}(s)$ (visible), and capability family set $\mathcal{F}$,
the \emph{Difference Matrix} $\text{DM}(s) = \langle\text{AD},\text{UD},\text{FC},\text{TN}\rangle$
partitions $\mathcal{F}$ into four disjoint cells:
\end{definition}

\vspace{-0.5em}
\begin{center}
\small
\begin{tabular}{@{}lcc@{}}
\toprule
 & $f \in \mathcal{C}_{\text{vis}}$ & $f \notin \mathcal{C}_{\text{vis}}$ \\
\midrule
$f \in \mathcal{C}_{\text{act}}$ & \textbf{AD} (aligned) & \textbf{UD} (hidden) \\
$f \notin \mathcal{C}_{\text{act}}$ & \textbf{FC} (false claim) & \textbf{TN} (absent) \\
\bottomrule
\end{tabular}
\end{center}

Summing over all families in $\mathcal{F}$ yields per-skill scalars
$\text{AD}(s)$, $\text{UD}(s)$, $\text{FC}(s)$, $\text{TN}(s)$.

\paragraph{Capability Recall (CR).}
What fraction of actual capabilities does the description disclose?
\begin{equation}
  \text{CR}(s) = \frac{\text{AD}(s)}{\text{AD}(s)+\text{UD}(s)}
  \label{eq:cr}
\end{equation}
$\text{CR}=1$ indicates perfect disclosure; $\text{CR}=0$ indicates total omission.
Prior work measured only CR~\cite{liao2020measuring,zhang2022measuring};
we show this is insufficient.

\paragraph{Description Precision (CP).}
What fraction of claimed capabilities are genuine?
\begin{equation}
  \text{CP}(s) = \frac{\text{AD}(s)}{\text{AD}(s)+\text{FC}(s)}
  \label{eq:cp}
\end{equation}
$\text{CP}=1$ means no false claims; defined only when $\mathcal{C}_{\text{vis}} \neq \varnothing$.
CP~$< 1$ identifies skills that promise capabilities developers never implemented.

\paragraph{Capability Discrepancy Index (CDI).}
The overall normalized inconsistency:
\begin{equation}
  \text{CDI}(s) = \frac{\text{UD}(s)+\text{FC}(s)}{\text{AD}(s)+\text{UD}(s)+\text{FC}(s)}
  \label{eq:cdi}
\end{equation}
CDI $\in [0,1]$; 0 = perfect consistency, 1 = complete mismatch.

\paragraph{Risk-Weighted Mismatch (RWM).}
CDI~treats all families equally; RWM weights undisclosed families
by their empirical malicious-injection risk:
\begin{equation}
  \text{RWM}(s) = \sum_{f \in \mathcal{C}_{\text{act}} \setminus \mathcal{C}_{\text{vis}}} w_f
  \label{eq:rwm}
\end{equation}
Weights $w_f$ are learned from T6 ground truth via logistic regression
(\S\ref{sec:groundtruth}).
RWM is used \emph{only for risk ranking}, not for IC classification---preserving
classification objectivity.

% ── IC Classification ────────────────────────────────────────────────────────
\subsection{Inconsistency Classification (IC-1--IC-4)}
\label{sec:ic_classification}

\begin{definition}[IC Classification]
\label{def:ic}
Given $\text{DM}(s) = \langle\text{AD},\text{UD},\text{FC},\text{TN}\rangle$,
skill $s$ is assigned to exactly one of six mutually exclusive, collectively
exhaustive classes: IC-0 (Accurate: $\text{CR}=\text{CP}=1$),
IC-1 (Full Omission: $\text{UD}>0, \text{AD}=0, \text{FC}=0$),
IC-2 (Partial Omission: $\text{UD}>0, \text{AD}>0, \text{FC}=0$),
IC-3 (False Claim: $\text{FC}>0, \text{UD}=0$),
IC-4 (Bidirectional: $\text{UD}>0, \text{FC}>0$),
or TN (No Capabilities: $\text{AD}=\text{UD}=\text{FC}=0$).
\end{definition}

The difference matrix induces a mutually exclusive, exhaustive
classification of every skill:

\begin{table}[h]
\centering
\caption{IC inconsistency types derived from the difference matrix.
  Prevalence over \clawhub{} ($N{=}18{,}443$, rule-based static extraction).}
\label{tab:ic_types}
\scriptsize
\begin{tabular}{@{}clp{3.2cm}rr@{}}
\toprule
\textbf{Type} & \textbf{Name} & \textbf{Condition} & \textbf{Count} & \textbf{\%} \\
\midrule
IC-1 & Full Omission      & UD$>$0, AD$=$0, FC$=$0 & 13{,}988 & 75.8\% \\
IC-2 & Partial Omission   & UD$>$0, AD$>$0, FC$=$0 & 2{,}513  & 13.6\% \\
IC-3 & False Claim        & FC$>$0, UD$=$0         & 78       &  0.4\% \\
IC-4 & Bidirectional      & UD$>$0, FC$>$0         & 538      &  2.9\% \\
\midrule
---    & Accurate           & UD$=$0, FC$=$0, AD$>$0 & 79       &  0.4\% \\
---    & No Capabilities    & AD$=$FC$=$UD$=$0       & 1{,}247  &  6.8\% \\
\midrule
\multicolumn{3}{l}{\textbf{Total Inconsistent (IC-1--IC-4)}} & \textbf{17{,}117} & \textbf{92.8\%} \\
\bottomrule
\end{tabular}
\end{table}

\paragraph{IC-1: Full Omission (CR\,=\,0).}
The description mentions \emph{no} capability from $\mathcal{C}_{\text{act}}$,
yielding CR\,=\,0.
This is the most prevalent pattern (75.8\%),
indicating that a large fraction of skills have marketplace descriptions
that provide no signal about detected executable capabilities.
Among T6 confirmed malicious skills, \textbf{22/34 (64.7\%)} are IC-1.
We note that rule-based extraction may over-attribute capabilities in
boilerplate code; IC-1 prevalence should be interpreted as an upper bound
pending LLM-based adjudication (future work).

\paragraph{IC-2: Partial Omission (0\,$<$\,CR\,$<$\,1).}
The description mentions some capabilities but omits others
(e.g., disclosing \textsc{File-Read}
while not mentioning \textsc{Cred-Read} or \textsc{Net-Write}).
This accounts for 13.6\% of skills;
\textbf{10/34 T6} skills (29.4\%) are IC-2.
Combined IC-1+IC-2 captures \textbf{94.1\%} of all T6 malicious skills.

\paragraph{IC-3: False Claim (CP\,$<$\,1).}
The description asserts a capability not detected in the code (0.4\%, $n{=}78$).
The most common false claims are
\textsc{Autonomous} and \textsc{Multi-Step}---skills marketed
as AI-orchestrated when no corresponding code pattern is detected.
No T6 skills are IC-3.

\paragraph{IC-4: Bidirectional Mismatch.}
Both omission and false-claim occur simultaneously (2.9\%, $n{=}538$):
the skill has unmentioned capabilities \emph{and} claims capabilities
not found in the code.
Example: ``raini-skill-audit'' has CRED\_READ+SHELL\_EXEC patterns
while claiming MULTI\_STEP without a corresponding implementation.

% ── Temporal Extension ──────────────────────────────────────────────────────
\subsection{Temporal IC: Version Update Patterns}
\label{sec:temporal_ic}

The difference matrix can be computed per version,
enabling IC-type tracking across updates.
We define three temporal patterns:

\paragraph{Temporal Capability Expansion (TCE).}
For transition $v_i \to v_{i+1}$:
\begin{equation}
  \text{TCE}(v_i \to v_{i+1}) = \text{RWM}(s, v_{i+1}) - \text{RWM}(s, v_i)
  \label{eq:tce}
\end{equation}
\emph{Silent expansion} occurs when TCE $> 0$ and the Semantic Disclosure Score
(SDS, changelog adequacy) $<0.35$.
Of \Ntransitions{} ClawHub version pairs,
\textbf{2,575 (10.6\%)} are silent expansions.

\paragraph{Silent Capability Injection (T1).}
When $\text{TCE} > 0$, $\text{SDS} < 0.35$, and new capabilities
include the credential-exfiltration pair
(CRED-READ $+$ NET-WRITE), we label the
transition a \emph{silent injection} (T1).
Of 512 T1 transitions, 100\% meet our thresholds.

\paragraph{T6: Confirmed Malicious Injection.}
Ground-truth malicious transitions ($n{=}40$) are identified via
developer-authored changelogs that explicitly acknowledge prior
malicious behavior (\S\ref{sec:groundtruth}).
T6 differs from T1 in verified intent rather than heuristic evidence.

% ── Detection Pipeline ──────────────────────────────────────────────────────
\subsection{IC-Based Detection Pipeline}
\label{sec:pipeline}

The IC framework drives a four-stage detection pipeline:

\begin{enumerate}[topsep=3pt,itemsep=2pt]
\item \textbf{Step 1 --- Capability Extraction}:
  SCS+DMR extracts $\mathcal{C}_{\text{act}}$ from all skill artifacts
  (SKILL.md, Python/JS/Shell source, dependency manifests, config files)
  using 300+ regex patterns across 14 LSCPS families; SCS alone
  matches code blocks within SKILL.md (89 family-core patterns);
  DMR extracts $\mathcal{C}_{\text{vis}}$ from marketplace description text.

\item \textbf{Step 2 --- Difference Matrix}:
  Compute AD, UD, FC, TN per family;
  derive CR, CP, CDI, RWM.

\item \textbf{Step 3 --- IC Classification}:
  Assign IC type (IC-1 through IC-4, Accurate, or TN-Only).
  SCS$\cup$DMR composite rank is the output of the core system.

\end{enumerate}

\noindent\textbf{Optional extensions.}
Steps 1--3 are the current core system: deterministic, $O(N)$ time,
no LLM required. Two optional extensions are designed for future deployment:
(a)~\textbf{Temporal Monitoring (TDM)}: invoked when version history is
available; detects silent capability expansions across updates;
(b)~\textbf{Semantic Adjudication (SA)}: a five-model LLM ensemble invoked
only for the top-ranked high-RWM subset (3--5\% of skills) to reduce
rule-based FPs and recover non-keyword expressed capabilities.
Neither extension participates in the current main evaluation;
both are reserved for future deployment refinement (\S\ref{sec:sa}, \S\ref{sec:temporal}).

% ── Evaluation Tasks ────────────────────────────────────────────────────────
\subsection{Research Questions}

\begin{itemize}[topsep=3pt,itemsep=2pt]
\item \textbf{RQ1 (Static IC):} What is the ecosystem-scale distribution
  of IC-1 through IC-4? Which capability families are most systematically
  hidden vs.\ falsely claimed? (\S\ref{sec:static})
\item \textbf{RQ2 (Temporal IC):} How does the IC type evolve across version
  updates? What fraction of updates constitute silent expansions? (\S\ref{sec:temporal})
\item \textbf{RQ3 (Detection):} Can the IC pipeline detect confirmed
  malicious supply-chain injections? What is the precision-recall tradeoff? (\S\ref{sec:groundtruth})
\item \textbf{RQ4 (Cross-platform):} Is capability-description inconsistency
  platform-specific or ecosystem-wide? (\S\ref{sec:static})
\end{itemize}
